\documentclass{article}
\usepackage[final]{neurips_2025}
\makeatletter\renewcommand{\@neuripsyear}{2026}\makeatother

\usepackage[utf8]{inputenc}
\usepackage[T1]{fontenc}
\usepackage{hyperref}
\usepackage{url}
\usepackage{booktabs}
\usepackage{amsmath}
\usepackage{graphicx}

\title{Dynamic Compute Allocation via Learned Halting for LLM Reasoning}

\author{Anonymous}

\begin{document}
\maketitle

\begin{abstract}
Test-time compute scaling improves LLM reasoning but exhibits systematic heterogeneity: 33.8\% of problems experience accuracy degradation with increased budgets (overthinking).
We propose Dynamic Halting, a learned controller that decides when to stop reasoning based on process dynamics rather than question features.
Our method achieves +9.6pp accuracy improvement over fixed baselines on GSM8K while using comparable compute.
\end{abstract}

\section{Introduction}
Scaling test-time compute via extended reasoning improves LLM performance, but we find this improvement is non-uniform.
Our analysis reveals 33.8\% of problems get worse with more compute (overthinking phenomenon).

We propose Dynamic Halting: a progressive reasoning framework with learned stopping criteria based on reasoning dynamics.

\section{Method}
See sections/method_dynamic.tex

\section{Experiments}
See sections/experiments_dynamic.tex

\end{document}
